\documentclass[12pt,journal,onecolumn]{IEEEtran}
\usepackage{array}
\usepackage{longtable}

\begin{document}

\title{Task 2: Architecture Framework}
\author{Team 27}
\maketitle

\section{Stakeholder Identification Based on IEEE 42010}

IEEE 42010 asks us to define stakeholders, their concerns and architecture views that answer those concerns.

\subsection{Stakeholders and Concerns}

\begin{longtable}{|p{3cm}|p{3cm}|p{8cm}|}
\hline
\textbf{Stakeholder} & \textbf{Role} & \textbf{Key Concerns} \\
\hline
Students & Main users for timed tests and practice & Simple test flow, fair anti-cheat handling, fast feedback, clear result summary \\
\hline
Administrators & Manage platform data and users & Reliable question/user management, monitoring visibility, clean role protection \\
\hline
Developers & Build and maintain prototype & Modular code boundaries, testability, clear APIs, low setup friction \\
\hline
Deployment Owner (Lab/Institution) & Runs service environment & Secure code execution, controlled infra cost, health visibility, recoverability \\
\hline
\end{longtable}

\subsection{Architecture Viewpoints and Views}

\subsubsection{Logical Viewpoint}
\begin{itemize}
    \item View: Internal module decomposition of one backend deployable unit.
    \item Main concerns: Maintainability, role isolation, feature evolution.
    \item Stakeholders: Developers, administrators.
    \item Summary: Backend is organized into modules: auth, users, questions, tests, submissions, analytics, monitor, evaluation, shared. Frontend is route-separated for student and admin flows.
\end{itemize}

\subsubsection{Process Viewpoint}
\begin{itemize}
    \item View: Runtime interaction flows for timed test and async submissions.
    \item Main concerns: Performance, responsiveness, reliability.
    \item Stakeholders: Students, developers.
    \item Summary: Timed-test flow is session-based and validated server-side. Async submission flow uses queue + worker. API returns quickly after enqueue, worker updates result later.
\end{itemize}

\subsubsection{Deployment Viewpoint}
\begin{itemize}
    \item View: Node API + worker processes with external dependencies.
    \item Main concerns: Setup simplicity, runtime isolation, operational clarity.
    \item Stakeholders: Deployment owner, developers.
    \item Summary: Backend API and worker are Node processes. Persistent data is in MongoDB (local Compass-compatible URI). Queue is Redis. Code execution isolation uses Docker runtime images per language.
\end{itemize}

\subsubsection{Security Viewpoint}
\begin{itemize}
    \item View: API access control + sandbox safety boundaries.
    \item Main concerns: Unauthorized access prevention, safe execution of untrusted code.
    \item Stakeholders: Students, admins, deployment owner.
    \item Summary: Global API auth gate protects all /api routes except signup/signin. Admin routes enforce role guard. Docker execution uses CPU, memory, timeout and network restrictions.
\end{itemize}

\subsubsection{Operational Viewpoint}
\begin{itemize}
    \item View: Health, logs, metrics and cache warm-up.
    \item Main concerns: Incident detection, debugging speed, read latency.
    \item Stakeholders: Developers, administrators.
    \item Summary: Structured request logs, monitor endpoints, worker telemetry, readiness probes and startup cache warm-up are used to improve operations and response-time consistency.
\end{itemize}

\section{Major Design Decisions (ADR Summary)}

This task has five accepted ADRs.

\subsection{ADR 001: Docker Containers for Code Execution}
\begin{itemize}
    \item Status: Accepted
    \item Why: Student code is untrusted and must not run in API process.
    \item Decision: Run code inside short-lived Docker containers with resource and network limits.
    \item Effect: Better safety and fault isolation, with extra execution overhead.
\end{itemize}

Full record: \texttt{ADRs/ADR-001-Docker-Code-Execution.md}

\subsection{ADR 002: Asynchronous Message Queue}
\begin{itemize}
    \item Status: Accepted
    \item Why: Evaluation work is slower than normal API requests.
    \item Decision: API enqueues submission jobs, worker processes jobs asynchronously.
    \item Effect: Better API responsiveness and burst handling, with added queue/worker complexity.
\end{itemize}

Full record: \texttt{ADRs/ADR-002-Async-Message-Queue.md}

\subsection{ADR 003: Modular Monolith}
\begin{itemize}
    \item Status: Accepted
    \item Why: Team size and timeline favor lower operational complexity.
    \item Decision: One deployable app with clear internal modules and strict boundaries.
    \item Effect: Easier debugging and onboarding, with less independent scaling than microservices.
\end{itemize}

Full record: \texttt{ADRs/ADR-003-Modular-Monolith.md}

\subsection{ADR 004: JWT Stateless Authentication}
\begin{itemize}
    \item Status: Accepted
    \item Why: Need role-aware auth without server-side session storage.
    \item Decision: Signed JWT with expiry and role claims; middleware-based checks per route.
    \item Effect: Fast auth checks and easy horizontal scale; early token revocation remains limited.
\end{itemize}

Full record: \texttt{ADRs/ADR-004-JWT-Authentication.md}

\subsection{ADR 005: MongoDB for Primary Data Storage}
\begin{itemize}
    \item Status: Accepted
    \item Why: Need flexible schema for mixed question types and nested execution outputs.
    \item Decision: Use MongoDB as primary document database with Mongoose validation.
    \item Effect: Faster schema evolution during prototype stage, with application-level integrity checks.
\end{itemize}

Full record: \texttt{ADRs/ADR-005-MongoDB-Database-Selection.md}

\end{document}
